\documentclass[aps,prd,reprint,nofootinbib]{revtex4-2}
% Co-theoretical-physicist derivation skeleton.
% Compile with pdflatex (revtex4-2). Swap prd -> prl/prb/pre as appropriate.

\usepackage{amsmath,amssymb,bm}
\usepackage{graphicx}

\begin{document}

\title{<Working title of the derivation>}
\author{<author>}
\affiliation{<affiliation>}
\date{\today}

\begin{abstract}
<One-paragraph statement of the result and the regime in which it holds.>
\end{abstract}

\maketitle

\section{Setup and assumptions}
% Model / action / Hamiltonian, boundary conditions, the regime, and the
% small (or large) parameter that organizes the expansion.

\section{Strategy}
% Name the method (perturbation theory, asymptotics, variational, self-consistency,
% symmetry, saddle point, ...) and justify why it suits this regime.

\section{Derivation}
% Step-by-step. Carry units. Track the order of every approximation.
\begin{equation}
  % key intermediate result
\end{equation}

\section{Result}
\begin{equation}
  \boxed{\; % final result \;}
\end{equation}

\section{Regime of validity and approximations}
% List each approximation as controlled (small parameter + estimate of neglected terms)
% or uncontrolled (flagged ansatz/truncation, to be tested numerically).

\section{Consistency checks}
% Limits recovered (free/classical/known), units, signs, symmetry, positivity/causality,
% and agreement with numerics (cross-reference the physics-numerics experiment).

\end{document}
